\documentclass{clean_cv}

% Add a BibTeX-style file encoding all of your publications to include here. You can export this from Zotero. Only include
% publications you want to appear here!
\addbibresource{publications.bib}
\author{José Luis Aguilar}
\headlineposition{Technology Consultant}

\begin{document}

\maketitle
% In this section, you can use any of the FontAwesome icons. The commands \faCenter and \faCenterStyle have been defined to properly center the icons
% when using the default font settings.
%
% You can use any of the icons listed in the fontawesome5 package documentation (https://ctan.math.utah.edu/ctan/tex-archive/fonts/fontawesome5/doc/fontawesome5.pdf)
% If you need to specify a specific style (as is done here for the address card), you should use the two-argument \faCenterCycle command
\begin{center}
\begin{tabular}{lll}
    \faCenter{envelope} \href{mailto:aguilarch.joseluis@gmail.com}{aguilarch.joseluis@gmail.com} & & \faCenter{phone-alt} +52 55 3468 2976 \\
    %\faCenter{orcid} %\href{https://orcid.org/0000-0002-1825-0097}{0000-0002-1825-0097} & 
    \faCenter{linkedin-in} \href{https://www.linkedin.com/in/josé-luis-aguilar-146b94194/}{José Luis Aguilar}& & \faCenter{github}  \href{https://github.com/JLAC11}{JLAC11}  %\faCenter{globe} \url{https://lotrproject.com} 
\end{tabular}
\end{center}
%Passionate student who loves being challenged. I believe no obstacle is too great, if goals and tools are appropriately selected. Bringing forth commitment to helping others and a sense of responsibility while being determined to impact on other people's lives. 

\vspace{-1em}

\section{Work Experience}
\begin{datetabular}{9em}
    \dateentry{Jan 2024 -- Ongoing}{
    \textbf{Technology Specialist $\rightarrow$ Technology Consultant - Qualtrics}
    {
    Currently maintaining and implementing brand health trackers for several clients in industries such as clothing, transportation, information security, software, and others. I am also responsible for building business logic, data transformations, dashboards, and analytics (Distinctive Image Associations, regression models, correspondence analysis); I maintain 9 accounts >10M USD of ACV, and have led 5 new implementations. I have been responsible for maintaining key automations within the Strategy \& Research team, as well as developing best practices for handling brand tracking datasets. Won Rookie of the Year award in 2024, and Innovation Award in 2025. I developed a new technique (brand slotting) that enables brand trackers to scale much better for multi-market studies.
    }
    }
\end{datetabular}

\begin{datetabular}{9em}
    \dateentry{Aug 2024 -- Ongoing}{
    \textbf{Consultant - Freelancing}
    {
    In my free time, I do freelancing projects. I have worked in banking, transportation, and laboratory/medical imaging sector projects, in commercial due diligences, headcount optimization, and financial projects, both in an IC and managerial role.
    }
    }
% There is something seriously funky happening between the tabular commands and the end of the itemize blocks.
% The command \eatvspace should be used if extra space appears at the end of a datetabular environment.
\end{datetabular}

\begin{datetabular}{9em}
    \dateentry{Jul 2022 - Jan 2024}{
    \textbf{Analyst - akya.}
    {I participated in a boutique strategy consulting firm. My collaborations were mainly in the financial sector by developing financial models, data science models (time series, survival analysis, regression,, and classification) and collecting data through several means (data scraping, pulling from several APIs, and creating small data processing workflows). Also participated in several projects for companies in the education, fitness and private equity sectors in Mexico.}
    }
% There is something seriously funky happening between the tabular commands and the end of the itemize blocks.
% The command \eatvspace should be used if extra space appears at the end of a datetabular environment.
\end{datetabular}
\vspace{-1.5em}

\section{Education}

% The datetabular environment takes one argument, which is the width of the left date column. As seen here:
%   9em is a good choice for "dual-date" formats (e.g. Sep 2015 - Nov 2019).
%   4em is a good choice for month/year dates (Sep 2014).
%   2em is a good choice for year-only dates (as seen in the publications)
\begin{datetabular}{9em}
% This is just a tabular environment, for the most part. The dateentry command has been defined for
% convienence. It takes two arguments, the first is the date and the second is whatever you wish placed to the right.
\dateentry{Aug 2018 -- Dec 2022}{
\textbf{Chemical Engineering -- Universidad Iberoamericana.}
{Won "Beca al Talento" award, which included full scholarship. Published a peer reviewed paper for a chemistry education journal. Took special courses on signal processing, machine learning, data science and analytics.}
}
\dateentry{Aug 2015 -- May 2018.}{
\textbf{High School -- Prepa Ibero.}
{Scholarship won by best score on admissions test. Worked in several projects focusing on the school community, such as developing a method to extract limonene from orange peels. Won \emph{Ser Prepa Ibero} medal, award given annually to a student embodying the values of solidarity, liberty, honesty, responsibility to society and contemplation.}
}
\end{datetabular}
\vspace{-1.5em}

\section{Publications}
\nocite{*} % Loads every entry from the attached .bib file
% Highlight takes three entries, given name, given name initials, and family name.
% If you have a middle initial, this call looks like:
% \highlightauthorname{Bob H.}{B. H.}{Smith}
\begin{datetabular}{9em}
%printbibyear has been defined to only print entries from a given citation year.
\dateentry{2021}{\printbibliography[heading=none]\vspace*{-2.5em}\ }
%\dateentry{3001}{\printbibyear{3011}}
\end{datetabular}

\section{Leadership and Teaching Experience}
\begin{datetabular}{9em}
    \dateentry{Aug 2017 -- Jul 2022}{
    \textbf{Private tutor, self employed, Mexico City.}{
     I work with students to develop personalized independent strategies, through strong listening skills and knowledge about diverse topics.}
    }
    \dateentry{Jan 2020 -- May 2020}{
    \textbf{Analytical Chemistry advisor -- Ibero Chemistry Department.}{ Required job for maintaining scholarship at college. Helped students develop strategies and gain insight on topics related to the subject.}    
    }
    \dateentry{Aug 2015 -- May 2018}{ 
    \textbf{Debate team member.} {Participated in ASOMEX debate tournaments. Won tournament at Tec de Monterrey (Nov. 2015) and in Jesuit schools in Acapulco, Puebla and Tlaxcala.}
    }
\end{datetabular}
\vspace{-1.5em}
\section{Skills \& Languages}
\begin{datetabular}{9em}
    \dateentry{Skills}{
        \textbf{Programming languages:} SQL, Python, MATLAB, C (beginner), Java (beginner), LaTeX, Javascript, R (programming language).
        
        \textbf{Software tools:} Simulink, Aspen Plus, Microsoft Suite, Google Cloud Platform, Qualtrics XM, QGIS, Jira, Git
        
        \textbf{Additional skills:} Communication skills, creative problem solving, statistical analysis.
    }
    \dateentry{Languages}{
    \textbf{Spanish:} Native speaker.
    
    \textbf{English:} Full proficiency.
    
    \textbf{French:} B1 certificate.
    }
\end{datetabular}

\end{document}